% !TeX program = lualatex
% !TeX encoding = utf8
% !TeX spellcheck = uk_UA
% !TeX root = ../test.tex

%=========================================================
\chapter{Хвилеводи}\label{\currfilebase}
\graphicspath{{\currfilebase/Pictures}}
%=========================================================

%% --------------------------------------------------------
\section{Від довгої лінії до хвилеводу}
%% --------------------------------------------------------

У попередньому розділі ми розглянули довгу лінію --- найпростішу напрямну систему, в якій електромагнітна хвиля поширюється вздовж \emphx{двох} провідників. Така конфігурація має принципове обмеження: зі зростанням частоти (а отже, зі зменшенням довжини хвилі) відстань між провідниками стає сумірною з довжиною хвилі, а втрати в діелектрику й металі різко зростають. На частотах вище кількох гігагерц коаксіальний кабель стає непридатним для передачі великої потужності.

Природний вихід --- \emphx{відмовитися від другого провідника}. Виявляється, що електромагнітна хвиля може поширюватися вздовж \emphx{порожнистої металевої труби} або \emphx{діелектричного стрижня}, не маючи жодного «зворотного» проводу. Такі напрямні системи називаються \emphx{хвилеводами}.

\begin{Attention}
	Хвилеводи --- це не екзотика. Прямокутні металеві хвилеводи використовуються в радарах, супутниковому зв'язку, НВЧ-печах. Оптичне волокно --- це діелектричний хвилевід, який складається з серцевини та оболонки з різними показниками заломлення. Саме завдяки волокну стала можливою сучасна глобальна мережа Інтернету: один оптичний хвилевід замінює тисячі мідних кабелів.
\end{Attention}

Будь-яку електромагнітну хвилю, що поширюється вздовж осі $z$, можна розкласти на поперечну (у площині $xy$) та поздовжню (уздовж $z$) складові. Залежно від того, які складові мають електричне й магнітне поля, розрізняють три типи хвиль:

\begin{itemize}
	\item \emphx{TEM-хвиля} (від англ. \emph{Transverse Electro-Magnetic}, \emphx{поперечно-електр\-омагнітна}): \emph{обидва} вектори --- і $\Efield$, і $\Hfield$ --- лежать у поперечній площині. Поздовжніх складових немає: $E_z = 0$, $H_z = 0$.
	\item \emphx{TE-хвиля} (від англ. \emph{Transverse Electric}, \emphx{поперечно-електрична}): електричне поле $\Efield$ лежить у поперечній площині ($E_z = 0$), але магнітне поле має поздовжню складову ($H_z \neq 0$). Іноді позначають як \emphx{H-хвиля}.
	\item \emphx{TM-хвиля} (від англ. \emph{Transverse Magnetic}, \emphx{поперечно-магнітна}): магнітне поле $\Hfield$ лежить у поперечній площині ($H_z = 0$), але електричне поле має поздовжню складову ($E_z \neq 0$). Іноді позначають як \emphx{E-хвиля}.
\end{itemize}

\begin{Attention}
	Літери в абревіатурах вказують на те, \emph{яке поле є повністю поперечним}:
	\begin{itemize}
		\item \emphx{TE} — \emph{E}лектричне поле поперечне (Transverse Electric);
		\item \emphx{TM} — \emph{M}агнітне поле поперечне (Transverse Magnetic);
		\item \emphx{TEM} — \emph{обидва} поля поперечні (Transverse Electro-Magnetic).
	\end{itemize}
\end{Attention}

У довгій лінії існує ТЕМ-хвиля, бо є \emph{два} провідники з різними потенціалами, між якими можна прикласти напругу й пропустити струм. У порожнистому металевому хвилеводі другого провідника немає, тому ТЕМ-хвиля в ньому існувати не може. Замість неї поширюються TE- або TM-хвилі.

\begin{Attention}
	Чому ТЕМ-хвиля неможлива в порожнистому хвилеводі? Інтуїтивно: для ТЕМ-хвилі електричне поле в перерізі має бути потенціальним, тобто описуватися скалярним потенціалом $\varphi(x,y)$, який задовольняє рівняння Лапласа $\Delta \varphi = 0$. Але граничні умови на стінках хвилеводу вимагають $\varphi = \text{const}$ на кожній стінці, а це разом із рівнянням Лапласа дає $\varphi = \text{const}$ всюди --- тобто $\Efield = 0$. Отже, ТЕМ-хвилі в порожнистому хвилеводі не буває.
\end{Attention}

%% --------------------------------------------------------
\section{Моди хвилеводу}
%% --------------------------------------------------------

Конфігурація поля в поперечному перерізі хвилеводу визначається рівняннями Максвелла разом із граничною умовою $\Efield_\parallel = 0$ на стінках. Розв'язки існують лише для дискретного набору пар цілих чисел $(m, n)$, кожна з яких задає \emphx{моду} хвилеводу. Індекси $m$ і $n$ (наприклад, $\text{TE}_{10}$) --- це \emphx{числа півхвиль} поля вздовж сторін хвилеводу.

Для прямокутного хвилеводу зі сторонами $a$ і $b$ граничні умови дають:
\[
	k_x = \frac{m\pi}{a}, \qquad k_y = \frac{n\pi}{b}.
\]
При цьому:
\begin{itemize}
	\item для \emphx{TE-мод}: $m, n \geq 0$, але не обидва одночасно нулі;
	\item для \emphx{TM-мод}: $m \geq 1$ і $n \geq 1$.
\end{itemize}

Чим більші $m$ і $n$, тим вища критична частота моди; основна мода має найменші індекси. Для хвилеводу з $a > b$ основна мода --- $\text{TE}_{10}$.

%% --------------------------------------------------------
\section{Прямокутний хвилевод: якісна картина}
%% --------------------------------------------------------

Найпростіший і найпоширеніший тип хвилеводу --- \emphx{прямокутна металева труба} зі сторонами $a$ і $b$ (зазвичай $a > b$). У ній можуть поширюватися хвилі типу $\text{TE}_{mn}$ або $\text{TM}_{mn}$.

Якісно картину можна уявити так: хвиля в хвилеводі --- це суперпозиція плоских хвиль, які багаторазово відбиваються від стінок. Щоб після відбиття хвиля «склалася сама з собою» (тобто утворила стійку картину), поперечні розміри хвилеводу мають бути кратні півдовжині хвилі. Саме тому моди дискретні, а не неперервні.

\begin{figure}[h!]
\centering
\begin{tikzpicture}[scale=1.1, >=latex]
	% --- Прямокутний хвилевід у перерізі ---
	\draw[thick] (0,0) rectangle (5,2.5);
	\node at (2.5,-0.4) {$a$};
	\node at (-0.4,1.25) {$b$};

	% --- Лінії електричного поля для TE10 ---
	\foreach \x in {0.6, 1.2, 1.8, 2.4, 3.0, 3.6, 4.2} {
		\draw[red!70!black, ->, thick] (\x, 0.2) -- (\x, 2.3);
	}
	\node[red!70!black] at (5.4, 2.0) {$\Efield$};

	% --- Лінії магнітного поля ---
	\draw[blue!70!black, thick, dashed] (0.3,1.25) .. controls (1.5,2.6) .. (2.5,1.25)
	                                    .. controls (3.5,-0.1) .. (4.7,1.25);
	\node[blue!70!black] at (5.4, 0.5) {$\Hfield$};

	% --- Підпис моди ---
	\node at (2.5, 3.1) {\emphx{Мода TE$_{10}$}};
\end{tikzpicture}
\caption{Структура поля найпростішої моди $\text{TE}_{10}$ у прямокутному хвилеводі. Електричне поле напрямлене поперек вузької стінки, а магнітне поле утворює замкнені петлі в поздовжньо-поперечній площині.}
\label{fig:te10_mode}
\end{figure}

Найпростіша мода --- $\text{TE}_{10}$. Вона має найнижчу критичну частоту, тому саме вона зазвичай використовується на практиці. Її структура показана на рис.~\ref{fig:te10_mode}: електричне поле напрямлене поперек хвилеводу (від однієї широкої стінки до іншої), а магнітне поле утворює замкнені петлі.

%% --------------------------------------------------------
\section{Критична довжина хвилі}
%% --------------------------------------------------------

Головна особливість хвилеводу: не будь-яка хвиля може в ньому поширюватися. Існує \emphx{критична довжина хвилі} $\lambda_{\text{кр}}$ (у вакуумі або в діелектрику хвилеводу), яка визначається розмірами хвилеводу. Для моди $\text{TE}_{10}$ у прямокутному хвилеводі:

\begin{equation}\label{eq:critical_wavelength}
	\tcbhighmath{
		\lambda_{\text{кр}} = 2a.
	}
\end{equation}

Якщо довжина хвилі генератора $\lambda_0 < \lambda_{\text{кр}}$, хвиля поширюється. Якщо $\lambda_0 > \lambda_{\text{кр}}$ --- хвиля \emphx{затухає} експоненційно і в хвилевод не проникає. Це фундаментальна відмінність від довгої лінії, де хвиля могла поширюватися на будь-якій частоті.

\begin{Attention}
	Фізичний зміст критичної довжини хвилі простий: щоб хвиля «вписалася» в хвилевід, її поперечний розмір має бути не меншим за півдовжини хвилі. Якщо $\lambda_0 > 2a$, хвиля просто «не поміщається» в хвилевід і відбивається назад. Саме тому хвилеводи --- це \emphx{фільтри верхніх частот}: вони пропускають високі частоти й блокують низькі. Це принципово відрізняє їх від довгої лінії, яка пропускає все.
\end{Attention}

Для довільної моди $\text{TE}_{mn}$ або $\text{TM}_{mn}$ критична довжина хвилі визначається формулою:

\begin{equation}\label{eq:critical_wavelength_general}
	\lambda_{\text{кр}} = \frac{2}{\sqrt{(m/a)^2 + (n/b)^2}}.
\end{equation}

Найменшу критичну довжину хвилі має мода з найменшими індексами. Для хвилеводу з $a > b$ це буде $\text{TE}_{10}$, для якої $\lambda_{\text{кр}} = 2a$. У діапазоні між $\lambda_{\text{кр}}^{(\text{TE}_{10})} = 2a$ і $\lambda_{\text{кр}}^{(\text{TE}_{20})} = a$ у хвилеводі поширюється лише одна мода --- це \emphx{одномодовий режим}, який найзручніший для практики. (Це справедливо, якщо $2b < a$; інакше наступною модою буде $\text{TE}_{01}$ з $\lambda_{\text{кр}} = 2b$.)

%% --------------------------------------------------------
\section{Фазова та групова швидкість у хвилеводі}
%% --------------------------------------------------------

У хвилеводі фазова швидкість хвилі виявляється \emphx{більшою} за швидкість світла у вакуумі:

\begin{equation}\label{eq:phase_velocity_waveguide}
	v_{\text{ф}} = \frac{c}{\sqrt{1 - (\lambda_0/\lambda_{\text{кр}})^2}} > c.
\end{equation}

Це не суперечить теорії відносності, бо фазова швидкість --- не швидкість передачі сигналу. Сигнал (енергія) передається з \emphx{груповою швидкістю}:

\begin{equation}\label{eq:group_velocity_waveguide}
	v_{\text{гр}} = c\sqrt{1 - (\lambda_0/\lambda_{\text{кр}})^2} < c.
\end{equation}

При наближенні $\lambda_0 \to \lambda_{\text{кр}}$ групова швидкість прямує до нуля: хвиля «зупиняється» і не переносить енергію. Це ще одна ілюстрація фізичного змісту критичної довжини хвилі.

\begin{Attention}
	Співвідношення $v_{\text{ф}} \cdot v_{\text{гр}} = c^2$ --- це загальна властивість хвиль у хвилеводах. Воно виглядає дивно, але насправді є прямим наслідком дисперсії: у хвилеводі фазова швидкість залежить від частоти, тому хвиля є дисперсною. У довгій лінії без втрат дисперсії немає, тому $v_{\text{ф}} = v_{\text{гр}} = c/\sqrt{\varepsilon_r}$.
\end{Attention}

%% --------------------------------------------------------
\section{Оптичне волокно --- діелектричний хвилевід}
%% --------------------------------------------------------

Металеві хвилеводи --- не єдиний тип напрямних систем. Якщо замість порожнистої труби взяти \emphx{діелектричний стрижень} із показником заломлення $n_1$, оточений середовищем із меншим показником $n_2 < n_1$, то електромагнітна хвиля може поширюватися вздовж стрижня завдяки \emphx{повному внутрішньому відбиттю}.

Саме так працює \emphx{оптичне волокно}. Його серцевина (core) має показник заломлення $n_1$, а оболонка (cladding) --- $n_2 < n_1$. Світло, яке входить у серцевину під достатньо малим кутом, багаторазово відбивається від межі «серцевина--оболонка» і поширюється вздовж волокна з дуже малими втратами.

\begin{Attention}
	Оптичне волокно --- це, по суті, хвилевід, у якому роль «стінок» відіграє не метал, а межа двох діелектриків. У ньому теж існують моди (їх називають $\text{LP}_{mn}$), теж є критична довжина хвилі (точніше, критична частота), і теж можливий одномодовий режим. Одномодове волокно з діаметром серцевини близько $9$ мкм використовується в магістральних лініях зв'язку --- саме воно забезпечує передачу сигналу на тисячі кілометрів без підсилення.
\end{Attention}

%% --------------------------------------------------------
\section{Чисельні приклади}
%% --------------------------------------------------------

Розглянемо типовий прямокутний хвилевід із розмірами $a = 2.3$ см, $b = 1.0$ см (це стандартний хвилевід X-діапазону, який використовується в радарах).

\paragraph{1. Критична довжина хвилі для моди $\text{TE}_{10}$.}
\[
	\lambda_{\text{кр}} = 2a = 4.6~\text{см}.
\]
Критична частота:
\[
	f_{\text{кр}} = \frac{c}{\lambda_{\text{кр}}} = \frac{3\times10^8}{0.046} \approx 6.5~\text{ГГц}.
\]
Отже, хвилевід пропускає частоти вище $6.5$ ГГц.

\paragraph{2. Робоча частота.}
Нехай генератор працює на $f = 10$ ГГц, тобто $\lambda_0 = 3$ см. Тоді:
\[
	\frac{\lambda_0}{\lambda_{\text{кр}}} = \frac{3}{4.6} \approx 0.652.
\]
Фазова швидкість:
\[
	v_{\text{ф}} = \frac{c}{\sqrt{1 - 0.652^2}} = \frac{c}{\sqrt{0.575}} \approx 1.32\,c \approx 3.95\times10^8~\text{м/с}.
\]
Групова швидкість:
\[
	v_{\text{гр}} = c\sqrt{0.575} \approx 0.758\,c \approx 2.27\times10^8~\text{м/с}.
\]
Перевірка: $v_{\text{ф}} \cdot v_{\text{гр}} = 1.32 \cdot 0.758\,c^2 \approx c^2$. \checkmark

\paragraph{3. Одномодовий режим.}
Наступна мода --- $\text{TE}_{20}$ з $\lambda_{\text{кр}} = a = 2.3$ см, тобто $f_{\text{кр}} \approx 13$ ГГц. Отже, в діапазоні $6.5$--$13$ ГГц у хвилеводі поширюється лише одна мода $\text{TE}_{10}$. Це і є робочий діапазон X-band радарів.

%% --------------------------------------------------------
\section{Задача}
%% --------------------------------------------------------

\begin{problem}
	Прямокутний хвилевід має розміри $a = 2.3$ см і $b = 1.0$ см.
	\begin{enumerate}
		\item Обчисліть критичну довжину хвилі та критичну частоту для мод $\text{TE}_{10}$, $\text{TE}_{20}$ і $\text{TE}_{01}$.
		\item Визначте діапазон частот, у якому хвилевід працює в одномодовому режимі.
		\item На частоті $f = 10$ ГГц обчисліть фазову та групову швидкості, а також довжину хвилі в хвилеводі.
		\item Поясніть, чому при наближенні частоти до критичної групова швидкість прямує до нуля.
	\end{enumerate}
\end{problem}

\begin{solution}
	\paragraph{1. Критичні довжини хвиль.}

	За формулою \eqref{eq:critical_wavelength_general}:
	\[
		\lambda_{\text{кр}} = \frac{2}{\sqrt{(m/a)^2 + (n/b)^2}}.
	\]

	Для $\text{TE}_{10}$ ($m=1$, $n=0$):
	\[
		\lambda_{\text{кр}} = 2a = 4.6~\text{см}, \qquad
		f_{\text{кр}} = \frac{c}{\lambda_{\text{кр}}} \approx 6.52~\text{ГГц}.
	\]

	Для $\text{TE}_{20}$ ($m=2$, $n=0$):
	\[
		\lambda_{\text{кр}} = a = 2.3~\text{см}, \qquad
		f_{\text{кр}} \approx 13.0~\text{ГГц}.
	\]

	Для $\text{TE}_{01}$ ($m=0$, $n=1$):
	\[
		\lambda_{\text{кр}} = 2b = 2.0~\text{см}, \qquad
		f_{\text{кр}} \approx 15.0~\text{ГГц}.
	\]

	\paragraph{2. Одномодовий режим.}

	У діапазоні від $6.52$ ГГц до найближчої вищої критичної частоти (тобто $13.0$ ГГц для $\text{TE}_{20}$) поширюється лише мода $\text{TE}_{10}$. Отже, одномодовий діапазон:
	\[
		6.52~\text{ГГц} < f < 13.0~\text{ГГц}.
	\]

	\paragraph{3. Фазова й групова швидкості на $f = 10$ ГГц.}

	Довжина хвилі у вакуумі:
	\[
		\lambda_0 = \frac{c}{f} = \frac{3\times10^8}{10^{10}} = 3~\text{см}.
	\]
	Відношення:
	\[
		\frac{\lambda_0}{\lambda_{\text{кр}}} = \frac{3}{4.6} \approx 0.652.
	\]
	Фазова швидкість:
	\[
		v_{\text{ф}} = \frac{c}{\sqrt{1 - 0.652^2}} \approx 1.32\,c \approx 3.95\times10^8~\text{м/с}.
	\]
	Групова швидкість:
	\[
		v_{\text{гр}} = c\sqrt{1 - 0.652^2} \approx 0.758\,c \approx 2.27\times10^8~\text{м/с}.
	\]
	Довжина хвилі в хвилеводі:
	\[
		\lambda_{\text{хв}} = \frac{v_{\text{ф}}}{f} = \frac{3.95\times10^8}{10^{10}} \approx 3.95~\text{см}.
	\]
	Цікаво, що $\lambda_{\text{хв}} > \lambda_0$ --- хвиля в хвилеводі «довша», ніж у вільному просторі. Це прямий наслідок того, що $v_{\text{ф}} > c$.

	\paragraph{4. Чому $v_{\text{гр}} \to 0$ при $f \to f_{\text{кр}}$.}

	При наближенні до критичної частоти $\lambda_0 \to \lambda_{\text{кр}}$, тому вираз під коренем прямує до нуля:
	\[
		v_{\text{гр}} = c\sqrt{1 - (\lambda_0/\lambda_{\text{кр}})^2} \to 0.
	\]
	Фізично це означає, що хвиля перестає переносити енергію вздовж хвилеводу: вона багаторазово відбивається від стінок під кутом, близьким до $90°$, і її поздовжня складова руху прямує до нуля. Енергія «застрягає» в поперечних коливаннях, і хвиля не поширюється.
\end{solution}